\subsection{Problem~1}%
\label{problem:1}

Find a condition on the numbers $a, b, c$ such that the following system of equations is consistent.
When that condition is satisfied, find all solutions
(in terms of $a, b, \text{and}, c$):
\begin{equation*}
  \systeme{x+3y+z=a,
  -x-2y+z=b,
  3x+7y-z=c}
\end{equation*}
Then choose the numbers $a, b, c$ to satisfy the conditions:
\begin{tasks}
  \task $x=0, y=3, z=1$
  \task $x=0, y=0, z=1$
  \task $x=-2, y=1, z=0$
\end{tasks}
For the computed numbers $a, b, c$, estimate the solutions using the Regularized FOCUSS algorithm.

\subsubsection*{Mathematics}

The system of independent linear equations is \textit{undetermined} if the
number of equations is greater than the number of unknown variables.
Such a system may be expressed in the form:
\begin{equation*}
  \matr{Ax}=\matr{b}, \quad \text{where} \quad
  \matr{A}\in\mathfrak{R}^{m\times{}n}, \quad
  \matr{b}\in\mathfrak{R}^m, \quad
  \matr{x}\in\mathfrak{R}^n, \quad \text{and} \quad m<{}n
\end{equation*}
where $\matr{x}_p$ is any particular solution and
$\matr{x}_h\in\mathcal{N}(\matr{A})$ is an arbitrary element of the null space
of $\matr{A}$ \citep{Meyer}.

\paragraph{Sparsity and the $\ell_p$ diversity measure.}
To single out one solution, additional prior knowledge is needed.
A common assumption is that the true solution is \textit{sparse} --- most of
its entries are zero \citep{Zdunek}.
This is formalised by minimising the $\ell_p$ quasi-norm diversity measure
introduced by \citet{GorodRao97}:
\begin{equation*}
  E^{(p)}(\matr{x}) = \sum_{j=1}^{n} |x_j|^p, \qquad 0 < p \leq 1,
\end{equation*}
subject to $\matr{Ax}=\matr{b}$.
For $p=1$ this is the convex LASSO penalty; as $p\to 0$ it approaches the
$\ell_0$ pseudo-norm (a count of non-zero entries).


\subsubsection*{Solution}

Following the steps of the \textit{Gauss-Jordan method}, we transform the matrix to the Row Reduced Echelon Form:
\begin{equation*}
\begin{split}
    \begin{amatrix}{1}{1}
        \matr{A} & \matr{b}
    \end{amatrix} =
    \begin{amatrix}{3}{1}
        \phantom{-}1 & \phantom{-}3 & \phantom{-}1 & a\\
        -1 & -2 & \phantom{-}1 & b\\
        \phantom{-}3 & \phantom{-}7 & -1 & c
    \end{amatrix} &\xrightarrow[R_3 + 3R_2]{R_2 + R_1}
    \begin{amatrix}{3}{1}
      \phantom{-}1 & \phantom{-}3 & \phantom{-}1 & a\\
      \phantom{-}0 & \phantom{-}1 & \phantom{-}2 & b + a\\
      \phantom{-}0 & \phantom{-}1 & \phantom{-}2 & c + 3b
  \end{amatrix} \xrightarrow[R_3 - R_2]{R_1 - 3R_2}\\
  \xrightarrow[R_3 - R_2]{R_1 - 3R_2}
  \begin{amatrix}{3}{1}
    \phantom{-}1 & \phantom{-}0 & -5 & -3b -2a\\
    \phantom{-}0 & \phantom{-}1 & \phantom{-}2 & b + a\\
    \phantom{-}0 & \phantom{-}0 & \phantom{-}0 & c + 2b - a
  \end{amatrix} &
\end{split}
\end{equation*}
Looking at the RREF we may easily conclude that for our system to be consistent:
\begin{equation*}
    c + 2b - a = 0
\end{equation*}
In such a case, $z$ is a free variable and the considered system has infinitely many solutions.

Now we may proceed with substituting the variables:
\begin{tasks}
    \task \begin{equation*}
      \text{for}\quad
  \begin{cases}
    x=0\\
    y=3\\
    z=1
  \end{cases}\quad
  \text{we get}\quad
  \begin{cases}
    a=0+3*3+1=10\\
    b=0-2*3+1=-5\\
    c=0+7*3-1=20
  \end{cases}
  \end{equation*}
  \task \begin{equation*}
  \text{for}\quad
  \begin{cases}
    x=0\\
    y=0\\
    z=1
  \end{cases}\quad
  \text{we get}\quad
  \begin{cases}
    a=0+3*0+1=1\\
    b=0-2*0+1=1\\
    c=3*0+7*0-1=-1
  \end{cases}
  \end{equation*}
  \task \begin{equation*}
  \text{for}\quad
  \begin{cases}
    x=-2\\
    y=1\\
    z=0
  \end{cases}\quad
  \text{we get}\quad
  \begin{cases}
    a=-2+3*1+0=1\\
    b=-(-2)-2*1+0=0\\
    c=3*(-2)+7*1-0=1
  \end{cases}
\end{equation*}
\end{tasks}

\subsubsection*{FOCUSS Estimation}

The Regularized FOCUSS (FOCal Undetermined System Solver) algorithm finds the
minimum $\ell_p$-norm solution over the affine solution set by iteratively
reweighting a Tikhonov-regularised least-squares problem.
Starting from a random initialisation $\matr{x}^{(0)}$, each iterate is

\begin{equation*}
  \matr{x}^{(k+1)}
  = \matr{W}_k^2 \matr{A}^T
    \!\left(\matr{A}\matr{W}_k^2\matr{A}^T + \lambda\matr{I}\right)^{-1}
    \matr{b},
  \qquad
  \matr{W}_k = \operatorname{diag}\!\left(\left|x_j^{(k)}\right|^{1-p/2}\right),
\end{equation*}

\noindent where $p\in(0,1]$ controls sparsity and $\lambda>0$ provides
regularisation.
Here we use $p=0.1$ and $\lambda=10^{-8}$.

\lstinputlisting[style=Matlab-editor]{problems/Problem_1.m}

The FOCUSS estimates for all three tasks are shown in Table~\ref{tab:focuss_p1}.

\begin{table}[h]
  \centering
  \begin{tabular}{ccc}
    \hline
    Task & True chosen $\matr{x}$ & FOCUSS estimate \\
    \hline
    (a) & $(0,3,1)^T$ & $(0,\;3,\;1)^T$ \\
    (b) & $(0,0,1)^T$ & $(0,\;0,\;1)^T$ \\
    (c) & $(-2,1,0)^T$ & $(0,\;\tfrac{1}{5},\;\tfrac{2}{5})^T$ \\
    \hline
  \end{tabular}
  \caption{FOCUSS estimates ($p=0.1$, $\lambda=10^{-8}$) for each task.}
  \label{tab:focuss_p1}
\end{table}
